\section{Simulation Strategy}
\label{sec:simulation-strategy}

Simulation in this thesis serves two specific purposes, and is not asked to do more. The first is geometric design-space filtering: before fabricating a physical configuration, a low-fidelity model is used to verify that the polyhedral arrangement, with its module placement and central skeleton, is mechanically reasonable, balance-free at rest, and capable of producing center-of-mass displacement under selective inflation. The second is control-policy benchmarking: a higher-fidelity rigid-body simulator is used to test candidate control rules across configurations before committing to physical experiments, which are slower and harder to reset between trials.

The two purposes are served by different tools. Low-fidelity geometric simulation is built in Blender (\textit{\cref{fig:sim_blender}}). Each polyhedral configuration is represented by spheres at the vertices of a chosen polyhedron, and the inflation of each module is approximated by scaling the sphere's radius. This level of fidelity is sufficient for a single question, namely whether a given polyhedral topology can shift its center of mass enough to tip over under selective inflation, and it returns an answer cheaply for any candidate polyhedron. Higher-fidelity control simulation is built in MuJoCo (\textit{\cref{fig:sim_mujoco}}). Modules are modelled as rigid spheres mounted on the central skeleton through ball-and-slide joint compositions, so each module can both rotate around its mounting point and translate radially as it inflates or deflates. Inflation is rendered as a slide-joint extension, which captures the geometric consequence of volume change but does not model the membrane's elasticity or the contact between adjacent inflated bodies.

This last gap is deliberate. The MuJoCo model gives an approximation of what the configuration would do if its modules were rigid extending balls rather than soft inflated membranes. Comparing the simulated motion against the physical robot's motion therefore isolates what the soft material itself is contributing, since everything except the material is held constant across the two. Sim-to-real differences are not errors in the model but measurements of the material's role, and the methodology treats the gap as evidence rather than as a problem to be eliminated.
